% !TeX root = ../main.tex
\section{Выводы}

В ходе выполнения лабораторной работы было проведено прогнозирование рабочего спутникового созвездия глобальных навигационных систем GPS, ГЛОНАСС, Galileo и BeiDou на $22$~ноября 2026~года ($15{:}27$ местного времени) для открытой локации в районе станции метро «Котельники» ($55{,}6741^\circ\text{ с.ш.}$, $37{,}8582^\circ\text{ в.д.}$), а также сопоставление расчётных параметров с результатами натурных измерений.

Установлено, что в заданный момент времени над горизонтом находился $41$ космический аппарат четырёх основных ГНСС при благоприятной пространственной геометрии ($HDOP = 0{,}32$, $VDOP = 0{,}51$). Устойчивый приём и включение в навигационное решение обеспечиваются для $31$ спутника ($75{,}6\%$ от прогнозируемого состава), расположенных на углах места свыше $10^\circ$. Показано, что энергетические характеристики принимаемых сигналов ($C/N_0$ в пределах от $20{,}8$ до $43{,}4\text{ дБ-Гц}$) напрямую определяются углом возвышения космического аппарата над горизонтом: максимальный уровень сигнала соответствует зенитной области, тогда как на малых углах места приём ограничивается атмосферным затуханием, диаграммой направленности антенны смартфона и частичным затенением горизонта элементами городской застройки.
